\Qhead{2.10}{Самовоздействие волновых пучков. Обращение фронта. Интенсивные акустические пучки}{Самовоздействие волновых пучков. Самофокусировка света. Приближения нелинейной квазиоптики и нелинейной геометрической оптики. Обращение волнового фронта. Интенсивные акустические пучки; параметрические излучатели звука.}
\Subfile{2.10.1}{sections/02_waves/q10/s01}
\Subfile{2.10.2}{sections/02_waves/q10/s02}
\Subfile{2.10.3}{sections/02_waves/q10/s03}

\recap{Самовоздействие меняет фазу и траектории самого пучка; дифракция препятствует сжатию. Фазосопряжение создаёт встречную волну с обращённой фазой. Акустическая квадратичная нелинейность позволяет формировать направленный разностный сигнал.}
\board{Нелинейное параболическое уравнение; $L_{\rm dif}$ и $L_{\rm nl}$; векторы четырёхволнового смешения; две акустические накачки и разностный пучок.}
\oral{Почему геометрическая оптика не может надёжно определить минимальный размер самофокусирующегося пучка?}
\bait{Можно оставить отличие фазосопрягающего зеркала от обычного: оно меняет знак пространственной фазы, необходимый для компенсации аберраций.}
\sources{Цысарь~\cite{tsysar2}, лекции~9--10, PDF-с.~37--46; университетские слайды по нелинейной оптике~\cite{nonlinearoptics}. Нормированное уравнение самовоздействия получено из параксиального уравнения; акустическое уравнение сопоставлено со стандартной формой KZK~\cite{kzk}.}
